\chapter{Introducción}

\section{Contexto y motivación}

El correo electrónico sobrevive a la proliferación de canales publicitarios por una razón económica: su coste por impacto es bajo y permite ajustar el mensaje al perfil de cada receptor. Esa ventaja, sin embargo, solo se materializa cuando la empresa acierta a quién contactar y con qué producto. Una campaña enviada de forma indiscriminada gasta el mismo presupuesto que una segmentada, pero rinde mucho menos; el problema, por tanto, no es enviar correos, sino decidir bien a quién dirigir cada uno.

En la práctica, los responsables de campaña se enfrentan a un problema de asignación de inversión con información incompleta: el comportamiento real del usuario ante un correo electrónico es difícil de anticipar, y los datos disponibles suelen ser escasos, ruidosos y con un fuerte desbalanceo entre interacciones positivas y negativas. La tasa real de clic en campañas de email marketing se sitúa habitualmente por debajo del 5\%, lo que convierte la predicción de comportamiento en un problema de clasificación con clases muy desbalanceadas.

El aprendizaje automático ofrece herramientas para abordar este problema de forma sistemática, combinando señales de comportamiento pasado con características del usuario y del producto anunciado. No obstante, su aplicación efectiva requiere decisiones metodológicas cuidadosas, especialmente en lo relativo al tratamiento del desbalanceo de clases, la escasez de interacciones observadas y la incorporación de información contextual externa.

\section{Pregunta de investigación y objetivos}

La pregunta central de este trabajo es la siguiente: \textit{¿en qué medida es posible predecir la probabilidad de interacción de un usuario con una campaña de email marketing a partir de su perfil demográfico, su historial de comportamiento y las características del producto anunciado, y qué tipo de señal resulta más determinante?}

Para dar respuesta a esta pregunta, se plantean los siguientes objetivos:

\begin{enumerate}
    \item Construir un pipeline de preprocesado que integre datos de campaña con información geográfica y demográfica de fuentes externas, generando perfiles de usuario completos y homogéneos.
    \item Desarrollar dos niveles de segmentación no supervisada: una segmentación demográfica pura para la asignación de usuarios nuevos (arranque en frío), y una segmentación comportamental basada en el historial de interacciones.
    \item Implementar un modelo de propensión que combine variables de usuario y de producto ---incluido el texto del asunto del correo---, comparando algoritmos y arquitecturas (plano frente a jerárquico), validado sin fuga de datos y con corrección explícita del sesgo de muestreo introducido por el entrenamiento balanceado.
    \item Construir un sistema de recomendación de productos y evaluarlo de forma controlada (partición temporal, \textit{baselines} de referencia y contraste estadístico), determinando qué señal (comportamental o demográfica) resulta más útil para el \textit{ranking}.
    \item Predecir el volumen semanal de clics mediante modelado de series temporales para apoyar la planificación de campañas.
\end{enumerate}

\medskip
\noindent\textbf{Enfoque y contribución.} El énfasis de este trabajo es metodológico. Más que maximizar una métrica de exactitud, se construye un pipeline reproducible de extremo a extremo (del dato bruto a un ejemplo de despliegue) que toma cada decisión de forma justificada: validación sin fuga de datos, comparación de cada modelo frente a \textit{baselines} triviales antes de adoptarlo, corrección del sesgo de muestreo y contraste estadístico de cada afirmación. La aportación no es un predictor de máximo rendimiento, sino un proceso que determina, con evidencia, qué señal es útil, qué modelo merece desplegarse y a qué coste. Por ello varias de las hipótesis siguientes se plantean precisamente para poder ser \emph{refutadas}: el valor del trabajo está tanto en confirmarlas como en descartarlas con fundamento. En suma, este TFM es una guía de cómo abordaría un equipo de datos un proyecto orientado a maximizar el ROI de las campañas: qué construir en cada etapa, cómo validarlo y cómo decidir qué merece desplegarse y qué no. El objeto de estudio es, por tanto, \emph{el método} tanto como los modelos: que algunos resultados sean negativos no resta valor al trabajo, sino que muestra el proceso funcionando como debe.

\medskip
\noindent\textbf{Hipótesis.} A partir de la pregunta y los objetivos, el trabajo contrasta cuatro hipótesis falsables:
\begin{itemize}
    \item \textbf{H1.} Las variables de \emph{producto} (sector, categoría y texto del asunto) mejoran de forma significativa la propensión al clic frente a un modelo que solo usa variables de usuario y frente al \textit{baseline} de popularidad.
    \item \textbf{H2.} Las variables demográficas enriquecidas aportan poder predictivo a la propensión al clic; se contrasta pese a su naturaleza sintética (Capítulo~\ref{cap:datos}), precisamente para medir cuánta señal individual conservan.
    \item \textbf{H3.} Un recomendador basado en el segmento comportamental del usuario supera al \textit{baseline} de popularidad en el \textit{ranking} de productos.
    \item \textbf{H4.} El modelo Prophet mejora la previsión del volumen de clics frente a \textit{baselines} ingenuos (naïve y media móvil).
\end{itemize}
El contraste de cada hipótesis se aborda, respectivamente, en los módulos M2 (Capítulo~\ref{cap:propension}; H1 y H2, esta última vía el estudio de ablación), M3 (Capítulo~\ref{cap:recomendador}; H3) y M4 (Capítulo~\ref{cap:forecast}; H4).

\section{Estructura de la memoria}

La memoria se organiza siguiendo el flujo del pipeline. El capítulo~\ref{cap:estado_arte} revisa el estado del arte en modelado de propensión, segmentación de clientes, sistemas de recomendación y predicción temporal aplicados al marketing digital. El capítulo~\ref{cap:datos} describe las fuentes de datos, el proceso de limpieza, el enriquecimiento demográfico-geográfico y el análisis exploratorio. El capítulo~\ref{cap:evaluacion} ofrece una visión general del pipeline y el marco de evaluación común (métricas, validación sin fuga y pruebas estadísticas). A continuación, cada módulo se presenta en su propio capítulo, exponiendo de forma conjunta su método y sus resultados: segmentación de clientes (Capítulo~\ref{cap:segmentacion}), modelo de propensión (Capítulo~\ref{cap:propension}), recomendador (Capítulo~\ref{cap:recomendador}) y previsión temporal (Capítulo~\ref{cap:forecast}). El capítulo~\ref{cap:negocio} cubre el despliegue de ejemplo y el análisis de retorno de la inversión. Finalmente, el capítulo~\ref{cap:conclusiones} recoge la discusión, las conclusiones y las líneas de trabajo futuro.
